\chapter{Translating a banking decision into a causal question}
\label{la:translation}

The preceding evidence review leaves a practical starting point: assistance
must preserve the question before it can improve an answer. This chapter
builds that question from population, action, outcome, loss and information
time. Incentive and servicing examples expose common substitutions; a
literature-graph study separates extracted claims from hypotheses; the
formulation workflow then leads to robust population comparisons and goal
clarification. Together these treatments explain which commitments an analyst
must resolve and which alternatives an assistant must retain.

\section{A business sentence can contain several questions}

``Should we increase the deposit offer?'' sounds like a single question, but
it leaves several choices unresolved. The intended action could be a change
in the quoted rate, the eligibility rule, the message shown to a customer or
the allocation of a fixed promotional budget. The outcome could be a new
deposit, retained balances, funding cost or total relationship value. The
population could be existing customers, applicants who received an offer or
everyone eligible under a proposed rule. Each choice defines a different
causal problem.

A language assistant can expose these alternatives and organize the evidence
needed to resolve them. It should not choose the easiest target to estimate
and silently treat that target as the owner's decision. In the project funding
case, the response of balances to a 25-basis-point increase and the change in
net funding cost over thirty days are retained as separate candidates. Their
business motivation is shared; their outcomes and decision consequences are
not identical. The unresolved choice is useful information, not a failure
that must be hidden by returning a number.

The co-pilot roadmap discussed by Alaa and colleagues motivates this progression
from a scientific question through assumptions, statistical targets and
interpretation. Their illustrative medical examples attend to eligibility,
baseline time, treatment and follow-up. The inspected paper also reports that
generic uncued answers can omit relevant design concerns. Applying those
principles to banking is a project proposal; the medical examples do not
validate banking clocks or interventions. \cite{v12_la_2407_19118}

\section{A question as a collection of commitments}

The project represents a question as
\begin{equation}
 Q=(\mathcal U,P,\mathcal A,\pi,\pi_0,Y,H,\nu,\ell,\mathcal I).
 \label{la:question}
\end{equation}
Here \(\mathcal U\) is the unit of analysis; \(P\) is the target population;
\(\mathcal A\) is the feasible action set; \(\pi\) and \(\pi_0\) are the
policy and comparator; \(Y\) is the outcome; \(H\) is the horizon; \(\nu\)
is the measure used to average over the population; \(\ell\) is the loss,
including its units; and \(\mathcal I\) specifies information available at
each action. This notation defines the project's declaration contract. It is
not an identification theorem.

The unit needs particular care. An account can be an observation while a
customer is the decision unit and a household is the unit within which actions
interact. An assistant can ask which role each identifier plays, but a column
called customer ID does not certify ownership or stable linkage. If several
accounts share a budget or pricing rule, an account-level action may also be
an infeasible description of the intervention. The feasible action set should
reflect what the responsible institution can actually change.

The comparator must be as explicit as the intervention. ``No treatment'' might
mean no new contact, continuation of the current contact policy, no rate
increase, or a different offer under a fixed budget. A dynamic policy can
include future actions conditional on history, so its comparator is also a
history-dependent rule. An effect relative to one comparator cannot be
reused for another solely because the same treatment variable appears in the
dataset.

The outcome and loss should also remain separate. A change in balances has
units of currency held over a horizon; funding expense combines balances,
rates and time. A larger deposit balance can have a positive acquisition
interpretation and an unfavorable net-cost implication. If \(g(a)\) is a
specified gross benefit and \(c(a)\) is a specified cost, the decision may
depend on \(E[g(a)-c(a)]\), not on \(E[g(a)]\) alone. The definitions and
valuation assumptions must be supplied before that subtraction is meaningful.
The assistant should make missing costs visible rather than invent prices.

\section{Changing a population measure changes the answer}

Let \(X\) denote baseline features and let
\(\tau(x)=E[Y(1)-Y(0)\mid X=x]\). A weighted target can be written
\begin{equation}
 \theta_w=\frac{E[w(X)\tau(X)]}{E[w(X)]},\qquad E[w(X)]>0.
 \label{la:weighted-target}
\end{equation}
Uniform, treated-population, overlap and trimmed targets use different
weights. Under the conditions that make a propensity \(e(x)=\Pr(A=1\mid X=x)\)
appropriate for the intended population, familiar choices are \(1\),
\(e(x)\), \(e(x)(1-e(x))\) and
\(1\{a\leq e(x)\leq b\}\), respectively. The formula records which average
is wanted. It does not establish exchangeability, support, estimation rates
or coverage.

An exact project example makes the distinction visible. Let two strata have
equal population probability, with propensities \(0.2\) and \(0.5\), and
effects \(0\) and \(10\). Trimming to \([0.1,0.9]\) keeps both strata, so
the trimmed-population effect is \(5\). Overlap weights are \(0.16\) and
\(0.25\), giving
\[
 \theta_{\mathrm{overlap}}=
 \frac{0.16(0)+0.25(10)}{0.16+0.25}=\frac{250}{41}
 =6.09756\ldots.
\]
Both calculations can be exact and still answer different questions. A label
such as ``overlap population'' is insufficient if its underlying weighting
rule is not retained.

This issue is directly relevant to the inspected OCI agent paper. Chou and
colleagues describe a principal/actor/critic workflow, reproducible notebooks,
point-in-time estimation and overlap diagnostics. The inspected trimming
passage uses a target label that must be reconciled with the actual weighting
rule before mapping it to conventional overlap weighting. The proposed
architecture is informative, but an assistant should not inherit a target
label without checking its definition. \cite{v12_la_2607_22443}

Changing the target to remove unsupported accounts can be scientifically
reasonable, but it is a new decision question. A bank owner may care about
precisely those accounts for which treatment support is weak. An honest report
can say that a narrower target is estimable under the supplied assumptions
while the original target remains unresolved. It should show the excluded
population, the new averaging rule and the implications for the decision.

\section{Information is governed by an availability clock}

The date on which an event occurred differs from the date on which an analyst
or policy could observe it. For a recorded variable \(V\), write \(e_V\) for
event time and \(a_V\) for availability time. The finite project fixtures
require \(e_V\leq a_V\), and a policy acting at time \(t\) can use the value
only if \(a_V\leq t\). A derived feature becomes available no earlier than
its required inputs. These rules prevent an apparently historical feature
from leaking future information into a decision.

In the servicing fixture, an account condition occurs at day ten but becomes
available at day eleven. A rule acting on day eleven may use it; a rule acting
on day ten may not. The distinction matters even when both values later appear
in the same warehouse table. A language assistant should request the recording
and availability convention, rather than infer it from a field name or a
sorted timestamp.

Advance announcements and forecasts require different variable definitions.
A forecast observed today of next month's condition is a current forecast
variable, not a measurement of a future realized condition with an impossible
availability time. Revised historical data similarly require a vintage or
revision convention. The project's finite validator does not implement every
revision protocol. It checks the stated recorded-value contract and leaves
more complex clocks as explicit extensions.

Equal timestamps do not determine direction. Two columns aggregated to the
same day may hide a finer action sequence or a common cause. The assistant
can propose a finer clock and identify the evidence needed to support it.
It cannot establish temporal precedence by placing one variable first in a
diagram. When the granularity is insufficient, the ambiguity should survive
into the model class.

\section{An available covariate is not automatically an adjustment variable}

In a sequential servicing model, \(A_0\) changes an intermediate condition
\(L_1\), which affects both a later action \(A_1\) and the outcome \(Y\).
The same \(L_1\) is a mediator for the first action and a common cause of the
second action--outcome relation. Its appropriate use depends on the target.
Including every available column in one regression does not preserve all
possible total, direct and policy effects.

The project declaration therefore distinguishes a policy-history purpose
from an ordinary adjustment purpose. A simple directed-path diagnostic flags
a post-treatment descendant proposed as an ordinary total-effect control.
It does not declare all post-treatment conditioning invalid. Controlled
direct effects and longitudinal procedures need their own question and
formal analysis. A warning from a declaration validator is not a replacement
for an adjustment criterion or an identification algorithm.

Observation processes can alter the target in a similar way. Suppose only
successfully linked records enter an analysis and both action and repayment
affect successful linkage. The selection indicator is part of the causal
problem, and an effect among linked records may differ from an effect in the
original population. A record-link score is an observed proxy; it is not
verified ownership. The observation fixture deliberately retains a latent
ownership variable, a selection node and an unresolved population choice.

CausalDS provides a useful literature distinction here: its conceptual SCM
variables and its separate observation model are not interchangeable. The
paper's narrated synthetic scenes can test how an agent handles that
distinction, but they do not supply authentic banking semantic or assignment
truth. \cite{v12_la_2607_08093}

\section{The handoff to formal analysis}

The assistant's useful output is a reviewable declaration and a record of
what remains unresolved. A selected declaration contains one explicit target,
its population measure, the intervention and comparator, clocks, feasible
actions, outcome and loss. An ambiguous declaration retains candidate targets
without covertly selecting one. A contradicted assumption cannot be left as
an inconsequential note while the calculation proceeds as if it held.

The local project validator exercises these distinctions with four synthetic
banking cases and adverse contract tests. It checks declared identities,
types, references and selected clock/graph conditions. It does not authenticate
an institution, certify a reviewer, determine whether a causal assumption is
true, or implement every MAG/PAG/CPDAG validity condition. A declaration ready
for formal analysis is ready for the next conditional operation, not a
qualified bank recommendation.

This boundary also explains the limited scope of the executable operation
tested with these declaration fixtures. A declaration can describe a dynamic
regime while that selected operation supports only a static binary population
ATE. The correct handoff then reports unsupported target mapping. Later
fitted-summary and servicing studies have their own declared regimes; their
existence does not qualify this particular mapping. Translating that request into a
static regression would produce a result at the cost of changing the user's
question. The appropriate remedy is to qualify a longitudinal operation or
obtain an explicit new target decision, with the original question preserved.
